problem:  
Let \((\mathbb{H}^n,\Delta,g_0,\mu)\) be the standard Heisenberg group endowed with its Carnot–Carathéodory distance \(d_{cc}\) and Haar measure \(\mu\), which satisfies the sharp measure contraction property \(\mathrm{MCP}(0,2n+3)\).  
Consider a smooth, *non–left-invariant* perturbation of the horizontal metric given by  
\[
\langle X,Y\rangle_{g_\varepsilon}(p) = \langle (I+\varepsilon\,A(p))X,Y\rangle_{g_0}, \qquad X,Y\in\Delta_p,
\]
where \(A(p)\) is a smooth, symmetric field of horizontal endomorphisms satisfying \(\operatorname{tr}A(p)=0\) and small \(C^2\)-norm \(\|A\|_{C^2}\ll1\). Let \(d_\varepsilon\) denote the induced Carnot–Carathéodory distance and \(\mu_\varepsilon\) the volume measure associated to \(g_\varepsilon\).  

**Question:**  
Does there exist a nontrivial \(A(p)\) (not identically zero and not arising from a left translation or automorphism of \(\mathbb{H}^n\)) such that \((\mathbb{H}^n,d_\varepsilon,\mu_\varepsilon)\) still satisfies an \(\mathrm{MCP}(0,N_\varepsilon)\) inequality with finite \(N_\varepsilon\)?  
If such perturbations exist, is \(N_\varepsilon\) necessarily greater than \(2n+3\), or can synthetic “nonnegative curvature” persist to first order in \(\varepsilon\)?  
Equivalently, can one identify an infinitesimal geometric invariant involving \(A(p)\), the horizontal curvature operator, and the Jacobian of the geodesic flow, that governs the stability (or instability) of the measure contraction property under non-homogeneous perturbations of the Heisenberg geometry?  

Why is it a "good" problem:  
This version breaks away from the known invariance regime—left-invariant anisotropic metrics—by probing *non-homogeneous*, spatially varying horizontal deformations. It asks whether the delicate balance that makes the Heisenberg group saturate \(\mathrm{MCP}(0,2n+3)\) is stable when the geometry is locally “tilted” in a non-symmetric fashion.  

It is “good” because:  
- **Novel territory:** No current theorem guarantees the stability or instability of the MCP for non-invariant sub-Riemannian perturbations of the Heisenberg metric. Understanding this small-\(\varepsilon\) behavior would forge new ground in the analysis of synthetic curvature conditions beyond homogeneous or model spaces.  
- **Conceptual simplicity:** The deformation \(g_\varepsilon\) is defined by an elementary geometric perturbation, yet its impact on MCP is deeply nontrivial and likely to require coupling tools from sub-Riemannian geodesic analysis, optimal transport, and measure theory.  
- **Interdisciplinary bridge:** It lies at the intersection of geometric analysis (horizontal curvature operators, geodesic expansion), synthetic Riemannian geometry (MCP/CD stability), and the PDE viewpoint under deformation of metric-measure structures.  
- **Potential impact:** A positive stability result would reveal that sub-Riemannian nonnegative synthetic curvature can survive mild inhomogeneity—revolutionizing assumptions of rigidity. A negative one would highlight an extreme sensitivity of Heisenberg’s synthetic curvature to local anisotropy, hinting at a new notion of *curvature spectrum* for Carnot-type spaces.  

The problem is thus minimal in statement, but resolving it demands uncovering the first-order analytic signature of synthetic curvature rigidity in the geometric heart of sub-Riemannian analysis.